\section{The spatial phase and all-time metric comparison}
\label{sec:phase-tracking}

Fix one graph from Theorem~\ref{thm:invariant-graph}, and
write $\phi_t$ for its complete finite flow. The phase
of a general initial metric will be defined by removing
this finite evolution from its parallel component.

\begin{theorem}\label{thm:all-time-tracking}
There are common $r_0,\delta,c,C>0$ and a $C^1$ map
$\Gamma_{p,w}:B_{X_p}(\delta)\to\mathfrak h$, with
uniformly bounded first derivative, having the following
properties. For every smooth initial metric in this ball,
let $u(t)$ be the actual flow and put
\[
 S(t)=\Pi u(t),\qquad
 y(t)=Qu(t)-\psi(S(t)),\qquad D_*=\|y(0)\|_X.
\]
Then
\begin{equation}\label{eq:phase-definition}
 \begin{split}
 \Gamma(u_0)&=\lim_{T\to\infty}\phi_{-T}(S(T)),\\
 S_{\rm eff}(t)&=\phi_t(\Gamma(u_0)),\qquad
 u_{\rm eff}(t)=S_{\rm eff}(t)+\psi(S_{\rm eff}(t))
 \end{split}
\end{equation}
are defined for all $t\ge0$. The reconstructed metric
$h_{\rm eff}=h_0e^{u_{\rm eff}}$ is an actual HYM solution
on the same fixed holomorphic bundle, and
\begin{equation}\label{eq:tracking-bounds}
 \|y(t)\|_X+\|u(t)-u_{\rm eff}(t)\|_X\le CD_*e^{-ct}.
\end{equation}
The centered relative logarithm
\begin{equation}\label{eq:centered-logarithm}
 \zeta(t)=\log\{e^{-u_{\rm eff}(t)/2}e^{u(t)}e^{-u_{\rm eff}(t)/2}\}
\end{equation}
satisfies $\|\zeta(t)\|_X\le CD_*e^{-ct}$, and hence
\begin{equation}\label{eq:metric-order-tracking}
 e^{-CD_*e^{-ct}}h_{\rm eff}(t)
       \le h(t)\le e^{CD_*e^{-ct}}h_{\rm eff}(t).
\end{equation}
All constants are independent of the base scales,
stabilizer type, and ratios of positive weak rates.
One can take $c=\sigma/2$, so the error is at most
$Cr^2D_*$ after $T_r=4\sigma^{-1}\log(1/r)$.

The phase equals $S_0$ on graph data $S_0+\psi(S_0)$.
It is unique among complete finite trajectories that
track $S(t)$ at an exponential rate greater than
$L_f=\|Df\|_\infty\le Cr^2$. It is jointly $K$-equivariant
and depends on the chosen graph and the full initial metric.
\end{theorem}

\subsection{Normal decay and the phase limit}

In the coordinates $S,y$, the extended equation is exactly
\begin{equation}\label{eq:normal-graph-equations}
 \dot S=f(S)+h(S,y),\qquad \dot y=-L_0y+F(S,y),
\end{equation}
where
\[
 \begin{split}
 h(S,y)&=\mathcal A(S,\psi(S)+y)-\mathcal A(S,\psi(S)),\\
 F(S,y)&=\mathcal B(S,\psi(S)+y)-\mathcal B(S,\psi(S))
                                                -D\psi(S)h(S,y).
 \end{split}
\]
The final term follows from differentiating the graph
coordinate and is essential. These maps vanish at $y=0$.
Their uniform derivatives satisfy
\begin{equation}\label{eq:normal-coupling-bounds}
 \|D_yh\|+\|D_yF\|\le\varepsilon,
 \qquad
 \|D_Sh(S,y)\|+\|D_SF(S,y)\|\le C\|y\|_X,
\end{equation}
where $\varepsilon$ is made small by choosing the cutoff
radius and then $r_0$. The second bound follows from
bounded mixed derivatives and vanishing at $y=0$.
For either sign of $t$, the finite variational equations give
\begin{equation}\label{eq:finite-backflow-bounds}
 \|D\phi_t\|\le e^{L_f|t|},\qquad
 \|D^2\phi_t\|\le K_f|t|e^{2L_f|t|},
 \qquad K_f=\|D^2f\|_\infty=O(r^2).
\end{equation}

Set $c=\sigma/2$. Choose the constants so that
$\varepsilon\int_0^\infty k(t)e^{ct}\,\dd t<1/2$.
The normal Duhamel formula then gives
\[
 \sup_{s\le T}e^{cs}\|y(s)\|_X
 \le MD_*+\tfrac12\sup_{s\le T}e^{cs}\|y(s)\|_X.
\]
It follows that $\|y(t)\|_X\le2MD_*e^{-ct}$.
This prevents normal exit for small initial $D_*$;
the bounded finite field and the integrable perturbation
$h$ prevent finite-time escape of $S$.

For $P_T=\phi_{-T}(S(T))$, differentiation cancels the
finite $f$ terms and gives
\begin{equation}\label{eq:phase-velocity}
        \frac{\dd P_T}{\dd T}
             =D\phi_{-T}(S(T))h(S(T),y(T)).
\end{equation}
Its norm is at most $C\varepsilon D_*e^{-(c-L_f)T}$.
Taking $L_f<c$ proves the phase limit and
$|\Gamma-S(0)|\le C\varepsilon D_*$.
For fixed $t$, apply the same differentiation to
$\phi_{t-T}(S(T))$ for $T\ge t$. Integration to infinity
gives
\begin{equation}\label{eq:finite-coordinate-tracking}
              |S(t)-S_{\rm eff}(t)|\le C\varepsilon D_*e^{-ct}.
\end{equation}
Together with the normal decay and $D\psi=O(r^2)$
this proves \eqref{eq:tracking-bounds} for the extended
system. If two finite trajectories track at rates larger
than $L_f$, applying $\phi_{-t}$ to their difference
forces their initial points to agree. Graph data have
$y=0$, so their phase is their initial finite coordinate.
This argument uses an upper expansion bound $L_f$;
it never divides by a positive weak rate.

\subsection{Differentiability of the phase}

The integrable smoothing kernel gives differentiable
dependence of mild solutions on each finite interval.
Let $W=\delta S$ and $Z=\delta y$. Their equations are
\[
 \begin{split}
 \dot W&=(Df+D_Sh)W+D_yh\,Z,\\
 Z(t)&=e^{-tL_0}Z(0)+\int_0^te^{-(t-s)L_0}
                     (D_SF\,W+D_yF\,Z)(s)\,\dd s.
 \end{split}
\]
Put $\gamma=\sigma/16$ and reduce $r_0$ so that
$L_f<\gamma/4$. With
$U_T=\sup_{t\le T}e^{-\gamma t}|W(t)|$ and
$V_T=\sup_{t\le T}e^{(c-\gamma)t}\|Z(t)\|_X$,
normal convolution and \eqref{eq:normal-coupling-bounds}
give
\[
 V_T\le C\|Z(0)\|_X+CD_*U_T,
 \qquad U_T\le C|W(0)|+C\varepsilon V_T.
\]
In the second inequality the coefficient $D_Sh$ is
integrable in time. For sufficiently small $D_*$,
absorption gives
\begin{equation}\label{eq:phase-variation-bounds}
 \begin{split}
 |W(t)|&\le Ce^{\gamma t}(|W(0)|+\varepsilon\|Z(0)\|_X),\\
 \|Z(t)\|_X&\le Ce^{-(c-\gamma)t}
                            (\|Z(0)\|_X+D_*|W(0)|).
 \end{split}
\end{equation}
The derivative of the velocity \eqref{eq:phase-velocity} is
\[
 D^2\phi_{-T}(S(T))[W(T),h]
      +D\phi_{-T}(S(T))(D_Sh\,W(T)+D_yh\,Z(T)).
\]
Its operator norm is bounded by a constant times
\[
 \varepsilon D_*(1+T)e^{-(c-2L_f-\gamma)T}
              +(D_*+\varepsilon)e^{-(c-L_f-\gamma)T}.
\]
This is uniformly integrable on a smaller initial ball.
The integral formula for $P_T$ therefore converges in
$C^1$, with continuous derivative by dominated convergence.
Composing with $S(0)=\Pi u_0$ and
$y(0)=Qu_0-\psi(\Pi u_0)$ proves the asserted phase
regularity. Equivariance follows from that of the
equations and uniqueness of the phase.

\subsection{Confinement in the actual equation}

Theorem~\ref{thm:spatial-entry} lets us reduce $\delta$
so that every smooth actual solution in the initial ball
satisfies $|S(t)|<a/4$ and $\|Qu(t)\|_X<a/4$ for all
positive time. Thus it remains where the extended and
actual equations agree. Local mild uniqueness identifies
the two solutions globally. Also
$D_*\le\|Qu_0\|_X+Cr^2|\Pi u_0|\le C\delta$.
Reduce $\delta$ again so that
\eqref{eq:finite-coordinate-tracking} is less than $a/4$.
Then $|S_{\rm eff}(t)|<a/2$, and the graph bound gives
$\|\psi(S_{\rm eff}(t))\|_X<a/4$. The selected graph
trajectory also stays in the actual region for all time.
Theorem~\ref{thm:invariant-graph} reconstructs its spatially
smooth actual HYM metric there.

The norm $X_p$ is a Banach algebra norm, since $D_0$ is
a derivation. Exponential and logarithm on a fixed small
positive matrix ball have bounded derivatives. Hence
\[
 (u,\widetilde u)\longmapsto
             \log(e^{-\widetilde u/2}e^ue^{-\widetilde u/2})
\]
is uniformly locally Lipschitz from $X_p\times X_p$ to
$X_p$ and vanishes on the diagonal. This proves the
bound for \eqref{eq:centered-logarithm}. Its matrix is
positive Hermitian with determinant one, so its spectral
bounds imply \eqref{eq:metric-order-tracking}. These
metric inequalities transport invariantly back to the
actual holomorphic bundle. The stated $C^1$ norm is the
one in the fixed transport used throughout the construction.
At $T_r$ we have $e^{-cT_r}=r^2$. This completes the
proof of Theorem~\ref{thm:all-time-tracking}.

\subsection{Finite preparation and all weak time scales}

The limit defining the phase has a finite preparation rule.
For any $T\ge0$, put $\Gamma_T=\phi_{-T}(\Pi u(T))$
and let $u_{{\rm eff},T}$ be the graph trajectory starting
at $\Gamma_T$. Integrating \eqref{eq:phase-velocity} gives
\begin{equation}\label{eq:finite-phase-error}
 |\Gamma_T-\Gamma|\le C\varepsilon D_*e^{-(c-L_f)T},
 \qquad |\Gamma_T-S(0)|\le C\varepsilon D_*.
\end{equation}
All these initial points can be put in a common graph
ball whose positive trajectories are actual and confined:
for small $S_1$, $\|S_1+\psi(S_1)\|_X\le C|S_1|$,
so Theorem~\ref{thm:spatial-entry} applies, and uniqueness
preserves the graph. A smaller initial $\delta$ puts
$\Gamma$ and every $\Gamma_T$ in that ball.

The two-solution distance inequality
\cite[Lemma 6.4.5]{Kobayashi}, applied to these two actual
graph flows, implies
\[
 \sup_x d_x(h_{{\rm eff},T}(t),h_{\rm eff}(t))
          \le\sup_x d_x(h_{{\rm eff},T}(0),h_{\rm eff}(0))
          \le C|\Gamma_T-\Gamma|.
\]
Here $d_x$ is the invariant distance on positive Hermitian
metrics in the fibre. On the common small ball it is
equivalent to the pointwise difference in logarithmic
coordinates. Apply the bounded projection $\Pi$ to
obtain $|S_T(t)-S_{\rm eff}(t)|\le C|\Gamma_T-\Gamma|$.
The graph derivative bound upgrades this finite-coordinate
bound to
\begin{equation}\label{eq:prepared-trajectory-error}
 \sup_{t\ge0}\|u_{{\rm eff},T}(t)-u_{\rm eff}(t)\|_X
                      \le C\varepsilon D_*e^{-(c-L_f)T}.
\end{equation}
The upgrade uses the finite graph; it does not assert
$C^1$ contraction for arbitrary metric flows.
Combining with the actual tracking estimate gives
\[
 \left\|\log\{e^{-u_{{\rm eff},T}(t)/2}e^{u(t)}
                           e^{-u_{{\rm eff},T}(t)/2}\}\right\|_X
       \le CD_*e^{-ct}+C\varepsilon D_*e^{-(c-L_f)T}.
\]
Since $L_f=O(r^2)$, taking $T=T_r$ makes the second
term $O(r^2D_*)$ uniformly for every later comparison time.
Moreover $T_r\lambda_j\to0$ for every positive weak rate.
Thus the preparation is negligible compared with all
inverse weak rates, including the times $h^{-2b}$ in
Proposition~\ref{prop:two-cycle-spectrum}.

The comparison is absolute tracking to the exact corrected
field $f$. It does not supply a relative error divided by
a vanishing trajectory or a nonlinear asymptotic exponent.
The initial phase retains spatial information and is
generally different from $\Pi u_0$. After shrinking the
initial basin so that its evolved points remain in the
phase domain, uniqueness also gives
$\Gamma(u(t))=\phi_t\Gamma(u_0)$.
Nonpolystable escape to a graded boundary is outside this
metric basin around an HE representative.
